\section{Introduction}
\label{sec:intro}

\begin{figure*}[t]
\centering
\resizebox{\textwidth}{!}{%
\begin{tikzpicture}[font=\scriptsize,
  blk/.style={draw=black!70,rounded corners=1pt,minimum height=3.2mm,inner sep=0pt},
  wait/.style={fill=red!25,draw=red!60,dashed},
  rel/.style={circle,draw=black!70,fill=white,inner sep=0.6pt,font=\tiny},
  sum/.style={rounded corners=3pt,draw,text width=#1,align=center,inner sep=3pt,font=\scriptsize}]
% ---------------------------------------------------------------- frame
\node[font=\small\bfseries] at (7.9,5.75) {When is the translation heard?};
\draw[dashed,gray!70,line width=0.6pt] (0,5.45)--(15.8,5.45);
\draw[dashed,gray!70,line width=0.6pt] (5.25,0.05)--(5.25,5.45);
\draw[dashed,gray!70,line width=0.6pt] (10.45,0.05)--(10.45,5.45);
\draw[dashed,gray!70,line width=0.6pt] (0,0.95)--(15.8,0.95);
\node[font=\footnotesize\bfseries,align=center] at (2.6,5.1) {1.\ Baseline: fixed rate};
\node[font=\footnotesize\bfseries,align=center] at (7.85,5.1) {2.\ Where the delay comes from};
\node[font=\footnotesize\bfseries,align=center] at (13.1,5.1) {3.\ Ours: playback-queue CRF};
% ---------------------------------------------------------------- shared timeline drawing (macro-free: two copies)
% panel 1: origin 0.45, 0.45 cm per second
\begin{scope}[shift={(0.45,0)}]
  \node[anchor=east,gray] at (0,4.35) {source A};
  \foreach \t/\h in {0.1/.12,0.3/.22,0.5/.16,0.7/.26,0.9/.1,1.1/.2,1.3/.25,1.5/.14,1.7/.22,1.9/.18,2.1/.27,2.3/.12,2.5/.2,2.7/.24,2.9/.15,3.1/.21,3.3/.26,3.5/.11,3.7/.19,3.9/.23,4.1/.17,4.3/.25,4.5/.13,4.7/.21,4.9/.27,5.1/.16,5.3/.2,5.5/.24,5.7/.12,5.9/.18}
    {\draw[spkc,line width=0.9pt] (0.45*\t,4.35-\h)--(0.45*\t,4.35+\h);}
  \node[anchor=east,gray] at (0,3.65) {text written};
  \foreach \i/\t in {1/2,2/4,3/6}{\node[rel] at (0.45*\t,3.65) {\i};}
  \node[anchor=east,gray] at (0,2.85) {heard};

  \node[blk,fill=frozen!30,minimum width=0.45*2.6cm,anchor=west] at (0.45*2,2.85) {1};
  \node[blk,fill=frozen!30,minimum width=0.45*2.6cm,anchor=west] at (0.45*4.6,2.85) {2};
  \node[blk,fill=frozen!30,minimum width=0.45*2.6cm,anchor=west] at (0.45*7.2,2.85) {3};
  \node[anchor=east,gray] at (0,2.3) {queue};
  \draw[wait] (0.45*4,2.18) rectangle (0.45*4.6,2.42);
  \draw[wait] (0.45*6,2.18) rectangle (0.45*7.2,2.42);
  \node[red!70!black,font=\tiny,anchor=west] at (0.45*7.3,2.3) {waits 1.2\,s};
  \draw[gray,densely dotted] (0.45*6,4.1)--(0.45*6,1.62);
  \draw[gray,densely dotted] (0.45*9.8,2.65)--(0.45*9.8,1.62);
  \draw[<->,>=Latex,red!70!black,thick] (0.45*6,1.72)--node[below,font=\tiny]{3.8\,s after A stops}(0.45*9.8,1.72);

\end{scope}
\node[sum=4.8cm,fill=red!12,draw=red!50] at (2.6,0.48) {Speech queues behind the speaker's own earlier speech: later blocks wait ever longer.};
% ---------------------------------------------------------------- panel 2: EndOffset decomposition (AMI Zh->En test)
\begin{scope}[shift={(6.0,0)},x=0.9cm]
  \node[anchor=east] at (-0.05,3.6) {B};
  \node[anchor=east] at (-0.05,2.5) {\sys{}};
  % B: model 0.67, detector 0.66, queue 0.65, last block 2.10 (seconds)
  \fill[gray!45] (0,3.45) rectangle (0.67,3.75); \fill[gray!20] (0.67,3.45) rectangle (1.33,3.75);
  \fill[red!45] (1.33,3.45) rectangle (1.98,3.75); \fill[frozen!40] (1.98,3.45) rectangle (4.08,3.75);
  \draw[black!60] (0,3.45) rectangle (4.08,3.75); \node[anchor=west] at (4.1,3.6) {4.08\,s};
  % MIST: 0.68, 0.66, 0.14, 1.64
  \fill[gray!45] (0,2.35) rectangle (0.68,2.65); \fill[gray!20] (0.68,2.35) rectangle (1.34,2.65);
  \fill[red!45] (1.34,2.35) rectangle (1.48,2.65); \fill[frozen!40] (1.48,2.35) rectangle (3.12,2.65);
  \draw[black!60] (0,2.35) rectangle (3.12,2.65); \node[anchor=west] at (3.14,2.5) {3.13\,s};
  \draw[decorate,decoration={brace,amplitude=3pt},thick] (1.33,3.82)--node[above=2pt]{after the text is written: 67\%}(4.08,3.82);
  \draw[decorate,decoration={brace,amplitude=3pt,mirror},thick] (0,3.38)--node[below=1pt,align=center,font=\tiny]{text timing\\(wait-$k$, HMT)}(0.67,3.38);
  \node[red!70!black,anchor=west,font=\tiny] at (1.55,2.2) {queue 0.65\,s $\rightarrow$ 0.14\,s};
  \foreach \c/\l/\x/\y in {gray!45/model lag/0.2/1.72,gray!20/detector wait/2.2/1.72,red!45/queue/0.2/1.47,frozen!40/last block/2.2/1.47}
    {\fill[\c] (\x,\y-0.08) rectangle (\x+0.22,\y+0.08); \node[anchor=west,font=\tiny] at (\x+0.25,\y) {\l};}
  \node[gray,font=\tiny\itshape] at (2.05,1.18) {EndOffset per unit, AMI Zh$\rightarrow$En test};
\end{scope}
\node[sum=4.8cm,fill=yellow!15,draw=orange!60] at (7.85,0.48) {Two thirds of the delay arise \emph{after} the text is written, where policies that time the text do not act.};
% ---------------------------------------------------------------- panel 3: playback CRF
\begin{scope}[shift={(11.15,0)}]
  \foreach \t/\h in {0.1/.12,0.3/.22,0.5/.16,0.7/.26,0.9/.1,1.1/.2,1.3/.25,1.5/.14,1.7/.22,1.9/.18,2.1/.27,2.3/.12,2.5/.2,2.7/.24,2.9/.15,3.1/.21,3.3/.26,3.5/.11,3.7/.19,3.9/.23,4.1/.17,4.3/.25,4.5/.13,4.7/.21,4.9/.27,5.1/.16,5.3/.2,5.5/.24,5.7/.12,5.9/.18}
    {\draw[spkc,line width=0.9pt] (0.45*\t,4.35-\h)--(0.45*\t,4.35+\h);}
  \foreach \i/\t in {1/2,2/4,3/6}{\node[rel] at (0.45*\t,3.65) {\i};}
  % rejected long rendering of block 1 and the chosen ones
  \node[blk,draw=gray,dashed,minimum width=0.45*2.6cm,minimum height=2.4mm,anchor=west,font=\tiny,text=gray] at (0.45*2,3.2) {long};
  \node[blk,fill=trained!35,minimum width=0.45*1.8cm,anchor=west] at (0.45*2,2.85) {1};
  \node[blk,fill=trained!35,minimum width=0.45*1.8cm,anchor=west] at (0.45*4,2.85) {2};
  \node[blk,fill=trained!35,minimum width=0.45*1.8cm,anchor=west] at (0.45*6,2.85) {3};
  \draw[->,>=Latex,densely dotted,thick,spkc!80!black] (0.45*6.9,2.62) to[bend left=28] node[below,font=\tiny,align=center]{backward message $\beta$: price of\\the wait a rendering causes later} (0.45*2.9,2.62);
  \draw[gray,densely dotted] (0.45*6,4.1)--(0.45*6,1.55);
  \draw[gray,densely dotted] (0.45*7.8,2.65)--(0.45*7.8,1.55);
  \draw[<->,>=Latex,spkc!70!black,thick] (0.45*6,1.62)--(0.45*7.8,1.62);
  \node[anchor=west,spkc!50!black,font=\tiny] at (0.45*7.9,1.62) {1.8\,s};
  \node[font=\tiny,align=left,anchor=west] at (0.45*8.3,3.55) {CRF state:\\backlog $h_i$\\$\times$ rendering $c_i$};
\end{scope}
\node[sum=4.8cm,fill=green!10,draw=green!50!black] at (13.1,0.48) {Choosing renderings by the waiting they cause later keeps the queue short: ASR-BLEU $\uparrow$, LAAL $\downarrow$.};
\end{tikzpicture}}
\caption{Motivation. (1)~A Talker that speaks every committed block in full at one rate builds a queue in each
speaker's stream, so later blocks are heard ever later. (2)~For the deadline baseline (B), two thirds of EndOffset arise
after the text is written, in the queue and the playback of the last block (Table~\ref{tab:decomp}); policies that
time the text (wait-$k$, MILk, HMT) act only before that. (3)~Playback alignment scores candidate renderings together
with the backlog in a linear-chain CRF whose backward messages charge each rendering the waiting it causes later
blocks (\S\ref{sec:method-plem}); \sys{} cuts the queue from 0.65 to 0.14\,s. Panels 1 and 3 are schematic; panel 2 is
measured on the AMI Zh$\rightarrow$En test set.}
\label{fig:motivation}
\end{figure*}


A meeting interpreter renders each contribution into the listener's language and makes clear whose it is: a
``yes'' from the project manager means something different from a ``yes'' from the client. And people talk at once:
over a third of our AMI test windows~\citep{carletta2005ami} hold several speakers. Simultaneous speech-to-speech
translation (SimulS2ST) systems translate one speaker and can keep that speaker's
voice~\citep{seamless2023,labiausse2025hibiki,cheng2025liveinterpret,deng2025simuls2s};
DiariST~\citep{yang2024diarist} and X-Translator~\citep{zhao2026xtranslator} attribute streaming \emph{text}
translation to speakers. We are not aware of a multi-speaker system that reports speech latency on overlapping
meetings.

Two gaps separate these systems from a meeting interpreter. \emph{Who said what}: during overlap, a speaker-blind
translator followed by speaker attribution keeps one speaker per span and loses the other voice. \emph{When it is
heard}: once speech is synthesised in each speaker's voice, a block can only play after that speaker's earlier speech,
so the listener's delay depends on how long the generated speech is, not only on when its text was written. For a
deadline-based streaming system, two thirds of the time from the end of an utterance to the end of its heard translation
arise after the text is written (Figure~\ref{fig:motivation}). Simultaneous policies time the
\emph{text}~\citep{ma2019stacl,arivazhagan2019milk,zhang2023hmt}; the playback of speech that is already written
has, to our knowledge, not been modelled.

We present \sys{} (\emph{Meeting Interpretation with Speaker Tracking}), in which one decoder writes, for every speaker
of a window, speaker-attributed, time-stamped translation segments $\spk{k}\,\tim{a}\,\text{translation}\,\tim{b}$,
which a speech generator speaks in the voice of speaker $k$ (Figure~\ref{fig:arch}). The decoder reads a frozen
streaming diarizer as speaker tokens interleaved with its acoustic tokens (adapted from G-STAR;~\citealp{peng2026gstar}),
so two concurrent voices stay two speakers, and writes on deadlines, each chunk one fixed lag after its source.

For the second gap we introduce \emph{playback alignment} (PA), a linear-chain CRF over a speaker's playback queue. Its
state is the backlog of synthesised speech and the rendering chosen for each block; deterministic queue dynamics make
the transitions, and the potentials combine the generator's likelihood, the fidelity of an independent recogniser and
the delay at which each word is heard. Because the backlog is a sufficient statistic of the past, forward--backward
sums exactly over all rendering sequences; a latency budget is imposed by posterior
regularisation~\citep{ganchev2010posterior}, whose dual is concave with a monotone derivative; and the posterior,
which prices each rendering by the queueing cost it imposes on all later blocks, is distilled into the speech generator.

Our contributions are (i) playback alignment, a CRF formulation of when and how fast committed translations are heard,
with exact inference, a posterior-regularised latency budget and posterior distillation (\S\ref{sec:method-plem});
(ii) speaker-first, time-stamped joint decoding with the diarizer read inside the decoder (\S\ref{sec:res-overlap});
and (iii) an online evaluation of whole meetings against rule-based playback layers given the same merging, selection
and even speaking rate, on three test sets and unused meetings (\S\ref{sec:res-pa}). Per utterance the
system remains closer to consecutive than to simultaneous interpretation; we therefore call it \emph{streaming}.
